% Preámbulo compartido (paquetes, encabezado, modo borrador).
% Los datos del proyecto/autor se definen en el main.tex de cada carpeta ANTES de incluir este archivo.
\usepackage[T1]{fontenc}
\usepackage[utf8]{inputenc}
\usepackage[spanish,es-noshorthands]{babel}
\usepackage[scaled=0.92]{helvet}
\renewcommand{\familydefault}{\sfdefault}
\usepackage{csquotes}
\usepackage{graphicx}
\usepackage{tabularx}
\usepackage{longtable}
\usepackage{enumitem}
\usepackage{float}
\usepackage{pdflscape}
\usepackage{tikz}
\usepackage{pgfgantt}
\usetikzlibrary{arrows.meta}
\usepackage{array}
\usepackage{fancyhdr}
\usepackage{lastpage}
\usepackage{titlesec}
\usepackage[margin=2.5cm,headheight=30pt]{geometry}
\usepackage[hidelinks]{hyperref}
\usepackage[style=apa,backend=biber]{biblatex}
\DeclareLanguageMapping{spanish}{spanish-apa}
\addbibresource{../master/referencias.bib}

\graphicspath{{../assets/}}
\newcommand{\logo}{\includegraphics[height=1cm]{uai-logo-redondo.jpg}}
\newcommand{\logoportada}{\includegraphics[height=3.5cm]{uai-logo-completo.jpg}}
\newcommand{\logofacultad}{\includegraphics[height=2.5cm]{uai-tecnologia-informatica-logo.png}}

% ---------- Niveles de numeración profundos (hasta 10.5.3.10.1) ----------
\setcounter{secnumdepth}{5}
\setcounter{tocdepth}{3}
\titleformat{\paragraph}[hang]{\normalfont\normalsize\bfseries}{\theparagraph}{1em}{}
\titleformat{\subparagraph}[hang]{\normalfont\normalsize\bfseries}{\thesubparagraph}{1em}{}
\titlespacing*{\paragraph}{0pt}{1.2ex plus 1ex minus .2ex}{0.8ex}
\titlespacing*{\subparagraph}{0pt}{1.2ex plus 1ex minus .2ex}{0.8ex}
\renewcommand{\theparagraph}{\thesubsubsection.\arabic{paragraph}}
\renewcommand{\thesubparagraph}{\theparagraph.\arabic{subparagraph}}

% ---------- Encabezado y pie de página ----------
\pagestyle{fancy}
\fancyhf{}
\lhead{\logo}
\chead{\nombreproyecto}
\rhead{\autor}
\cfoot{Página \thepage{} de \pageref{LastPage}}
\renewcommand{\headrulewidth}{0.4pt}
\fancypagestyle{plain}{\fancyhf{}\cfoot{Página \thepage{} de \pageref{LastPage}}}

% ---------- Modo borrador ----------
% Inserta un marcador bajo cada encabezado para que las secciones vacías
% puedan paginarse. Poner \borradorfalse al terminar la redacción.
\newif\ifborrador
\newif\ifcapredactado
\borradortrue

% ---------- Entregable ----------
% \tfifalse: Seminario (Cap. 1-5, 2 CU). \tfitrue: Trabajo Final (Cap. 1-9, 4 CU).
\newif\iftfi
\tfifalse
% \parcialtrue: entrega parcial (hasta el Cap. 9, sin carpeta técnica, contenido ni anexos).
\newif\ifparcial
\parcialfalse
% Tipo de documento mostrado en la carátula.
\newcommand{\tipodoc}{Seminario de Trabajo Final}
\ifborrador
  % Un capítulo ya redactado declara \capredactadotrue al inicio de su texto.tex
  % (el marcador se omite); cuerpo.tex lo reinicia antes de cada capítulo.
  \newcommand{\pend}{\ifcapredactado\else\par\textit{[Pendiente de redacción]}\par\fi}
  \let\oldsubsection\subsection
  \let\oldsubsubsection\subsubsection
  \let\oldparagraph\paragraph
  \let\oldsubparagraph\subparagraph
  \renewcommand{\subsection}[1]{\oldsubsection{#1}\pend}
  \renewcommand{\subsubsection}[1]{\oldsubsubsection{#1}\pend}
  \renewcommand{\paragraph}[1]{\oldparagraph{#1}\pend}
  \renewcommand{\subparagraph}[1]{\oldsubparagraph{#1}\pend}
\fi

